% SPDX-License-Identifier: Apache-2.0
\section{A Value Wider Than 128 Bits}
\label{sec:x-wide}

A \code{U<N>} is at most 128 bits, one Rust integer. A wider value
names how many 128-bit limbs it takes: \code{U<256, 2>} is 256 bits
in two. The unit is an accumulator at that width, cleared on request,
which adds its input each cycle and puts out the sum and two bytes
of it, taken with \code{slice}: the byte just above the low limb,
where a carry lands, and the top byte. The run fills the low limb, adds
one so the sum carries into the high limb, sets the top bit, and then
adds all ones, which is minus one and takes one away. The netlists
declare 256-bit vectors, and both are checked under nvc and Verilator
against the trace, as Section~\ref{sec:x-checked} describes, so the
trace, the testbenches and the two languages are all shown to hold
the full width.

\begin{figure}[htbp]
\centering
\input{wide_timing.tex}
\caption{The 256-bit accumulator, drawn as two bytes of its sum: the
carry into the high limb, then the top bit. The full values are in
the printout.}
\label{fig:x-wide}
\end{figure}

\lstinputlisting[caption={\code{ex\_wide.rs}. Run with \code{bazel run //lib/examples:ex\_wide}.},label={lst:x-wide}]{lib/examples/ex_wide.rs}

\lstinputlisting[style=ascii,caption={What \code{ex\_wide} printed when the build ran it: the sums as two limbs, then the lowering.},label={lst:x-wide-out}]{out_wide.txt}
